% Análisis estratégico y encaje formativo derivados del informe complementario.
\chapter{Análisis DAFO y justificación}
\label{cap:dafo}

\section{Análisis DAFO}
\label{sec:dafo}

El análisis se refiere al proyecto académico y a su contexto, no a un producto
ya desplegado. Los factores de mercado proceden de la revisión documental del
capítulo~\ref{cap:estado-arte}; las capacidades internas son hipótesis de
diseño que deberán contrastarse durante el desarrollo \citep{votedirect2026comparativa,specter2020voatz,haines2020swisspost}.

\small
\begin{longtable}{|p{0.46\textwidth}|p{0.46\textwidth}|}
\caption{Matriz DAFO de la propuesta de plataforma.}\label{tab:dafo}\\
\hline
\textbf{Debilidades (internas)} & \textbf{Fortalezas (internas)} \\
\hline
\endfirsthead
\hline
\textbf{Debilidades (internas)} & \textbf{Fortalezas (internas)} \\
\hline
\endhead
\hline
\endfoot
\hline
\endlastfoot
Equipo unipersonal y tiempo limitado por el TFM; ausencia inicial de certificaciones, historial electoral y auditoría independiente; coste de operar infraestructura disponible y segura; riesgo de defectos si se implementa criptografía propia. & Arquitectura cloud multi-tenant y API-first como objetivo; posibilidad de código abierto y transparencia; libertad para diseñar una UX actual; acceso potencial a asociaciones y grupos universitarios para validar necesidades. \\
\hline
\textbf{Amenazas (externas)} & \textbf{Oportunidades (externas)} \\
\hline
Competidores con ofertas gratuitas o de bajo coste; desconfianza hacia el voto remoto por fallos públicos documentados; obligaciones del RGPD y estatutos internos; phishing, malware, DDoS y ciclos de compra institucional prolongados. & Integración SSO, API y webhooks poco uniformes según la revisión; necesidad de PWA accesible; oportunidad de recibos y auditoría explicables; oportunidad de precios claros, actas exportables y localización al mercado hispanohablante. \\
\hline
\end{longtable}
\normalsize

\section{Estrategias que se desprenden}
\label{sec:estrategias-dafo}

\begin{itemize}
  \item \textbf{Fortalezas y oportunidades:} orientar el diseño a API REST
  documentada, autenticación federada opcional y PWA accesible, publicando decisiones y
  evidencias de auditoría.
  \item \textbf{Debilidades y oportunidades:} evitar criptografía propia y
  concentrar el trabajo en integración, aislamiento, experiencia de uso y
  pruebas reproducibles; incorporar componentes criptográficos revisados solo
  si el alcance y la revisión especializada lo permiten.
  \item \textbf{Fortalezas y amenazas:} responder a la desconfianza mediante
  transparencia de requisitos, código y límites conocidos, además de pruebas
  documentadas.
  \item \textbf{Debilidades y amenazas:} limitar la primera versión a
  consultas y votaciones internas de asociaciones u organizaciones, sin
  presentarla como apta para elecciones oficiales vinculantes.
\end{itemize}

\section{Encaje con el máster}
\label{sec:encaje-master}

% [Bloque B] Experiencia profesional, Prácticas Externas y asignaturas.
El trabajo se apoya en las competencias del máster, pero también en la
experiencia profesional y en las Prácticas Externas del autor, que han
condicionado las decisiones de arquitectura. Tres años de desarrollo full
stack en el sector bancario, centrados en microservicios Java/Spring Boot con
arquitectura hexagonal, mensajería con Apache Kafka y publicación de APIs
mediante Apigee, aportaron familiaridad con sistemas distribuidos en los que
la separación de responsabilidades, los contratos de interfaz y la
trazabilidad son requisitos de negocio. Esa experiencia motivó que VeriVoto se
diseñara como un conjunto de cinco servicios con base de datos propia y
aislamiento por organización (capítulo~\ref{cap:diseno-tecnico}), en lugar de
un monolito, y que la urna se separase físicamente del censo.

Las Prácticas Externas en Noa Estudio (Noa Software, S.~L.; 250 horas entre
septiembre y octubre de 2026) complementaron ese perfil, centrado sobre todo en
el backend, en la capa de cliente y en la seguridad de aplicaciones web. En
TaskFlow se desarrolló una API REST con autenticación JWT, contraseñas
protegidas con bcrypt y control de acceso por roles, junto con un cliente web
en Angular (guards de ruta e interceptores HTTP) y una aplicación móvil en
Ionic con Capacitor. En TaskVoid, un producto en uso, se implementaron
notificaciones Web Push con identificación del servidor mediante VAPID y un
Service Worker, y se trabajó con contenedores Docker, nginx, integración
continua en GitHub Actions y Pull Requests revisados por el equipo. Esta
práctica proporcionó la soltura técnica necesaria para construir la PWA
Angular/Ionic de VeriVoto y su modelo de autorización basado en tokens JWT y
roles por organización.

La tabla~\ref{tab:encaje-master} relaciona cada asignatura con su aportación
al proyecto y con el contexto en el que esa competencia se aplicó antes de
este trabajo.

\small
\begin{longtable}{|>{\RaggedRight\arraybackslash}p{0.18\textwidth}|>{\RaggedRight\arraybackslash}p{0.33\textwidth}|>{\RaggedRight\arraybackslash}p{0.39\textwidth}|}
\caption[Relación entre asignaturas del máster, proyecto y aplicación previa.]{Relación entre asignaturas del máster, aportaciones al proyecto y aplicación previa en prácticas o experiencia profesional.}\label{tab:encaje-master}\\
\hline
\textbf{Asignatura} & \textbf{Aportación al proyecto} & \textbf{Aplicación en prácticas/experiencia} \\
\hline
\endfirsthead
\hline
\textbf{Asignatura} & \textbf{Aportación al proyecto} & \textbf{Aplicación en prácticas/experiencia} \\
\hline
\endhead
\hline
\endfoot
\hline
\endlastfoot
Cloud Computing & Diseño del despliegue SaaS multi-tenant, aislamiento, escalabilidad y operación en nube. & Prácticas: contenedores Docker y \texttt{docker-compose} para levantar base de datos, backend y frontend; nginx y configuración mediante variables de entorno. La plataforma de despliegue de la empresa era OpenShift, a la que no se tuvo acceso (sección~\ref{sec:limitaciones-trabajo-futuro}). \\
\hline
Desarrollo \mbox{Web/Móvil} & Interfaz adaptable, PWA y API REST. & Prácticas: cliente Angular con guards, interceptores y panel en TaskFlow; aplicación Ionic/Capacitor; Service Worker y Web Push en TaskVoid. Experiencia: APIs REST con Spring Boot en el sector bancario. \\
\hline
Seguridad Informática & Autenticación, modelado de amenazas, privacidad, pruebas y auditoría. & Prácticas: autenticación JWT, contraseñas con bcrypt y control por roles en TaskFlow; claves VAPID como variables de entorno y suscripciones push restringidas por rol; actualización de dependencias con \texttt{npm audit}. \\
\hline
Integración de Tecnologías y Servicios & Proveedor de identidad opcional, API, webhooks y servicios de notificación. & Experiencia: APIs de proceso en Spring Boot expuestas mediante Apigee e integración asíncrona con Apache Kafka. Prácticas: servicio de notificaciones Web Push, Pull Requests e integración continua con GitHub Actions. \\
\hline
Bases de Datos a Gran Escala & Modelo multi-tenant, separación de datos y registro de auditoría. & Prácticas: modelado de datos con Prisma sobre MySQL/MariaDB; relaciones entre clientes, proyectos y tareas, y trazabilidad mediante historial en TaskVoid. \\
\hline
Visualización de Datos & Paneles de participación y presentación de resultados agregados. & Prácticas: panel con gráficas (ng2-charts) en TaskFlow y módulo de Planificación con diagrama de Gantt en TaskVoid. \\
\hline
Planificación de las TIC & Planificación, estimación de costes y análisis de alternativas de servicio. & Prácticas: planificación semanal, estimación en horas y Pull Requests con plantilla de pruebas y despliegue. Experiencia: gestión ágil con Jira y Trello y control de versiones con Git Flow. \\
\hline
Inteligencia Artificial / Sistemas Inteligentes & Extensión opcional futura para estudiar anomalías en patrones de participación, sin inferir opciones de voto. & Prácticas: uso de un asistente de IA en el flujo de desarrollo y en la elaboración de planes de prueba; no se aplicaron técnicas de IA al producto. \\
\hline
\end{longtable}
\normalsize

Las prácticas pusieron de manifiesto, además, una carencia que ha
condicionado este trabajo. La validación se apoyaba en pruebas manuales
documentadas en cada Pull Request, y la revisión del equipo detectaba
defectos en casos límite que esas pruebas no cubrían. En VeriVoto esta
lección se traduce en pruebas automatizadas desde el dominio (unitarias, de
arquitectura y de integración con PostgreSQL) y en análisis estático y de
dependencias en la integración continua, aunque su cobertura sigue siendo
parcial (capítulo~\ref{cap:evaluacion-resultados}).

\section{Delimitación del MVP y trabajo futuro}
\label{sec:alcance-mvp}

El catálogo detallado de la memoria conserva los códigos RF-01 a RF-19 ya
usados por los casos de uso. El informe complementario propone otra numeración
(RF-01 a RF-30); la columna «Origen / hueco» identifica esos códigos de origen
sin reasignar los existentes. La tabla~\ref{tab:alcance-mvp} explicita la
selección de alcance y permite revisar cada decisión con el tutor.

\small
\begin{longtable}{|p{0.23\textwidth}|p{0.27\textwidth}|p{0.39\textwidth}|}
\caption{Decisión provisional del alcance de requisitos funcionales.}\label{tab:alcance-mvp}\\
\hline
\textbf{Decisión} & \textbf{IDs locales / origen} & \textbf{Criterio de alcance} \\
\hline
\endfirsthead
\hline
\textbf{Decisión} & \textbf{IDs locales / origen} & \textbf{Criterio de alcance} \\
\hline
\endhead
\hline
\endfoot
\hline
\endlastfoot
MVP propuesto & RF-01--10, RF-12--19 y RF-21 & Se cubren mediante requisitos funcionales y no funcionales del catálogo local; el detalle de implementación y validación queda sujeto a planificación y aceptación del tutor. \\
\hline
MVP acotado & RF-20 (informe RF-11); RF-21 (informe RF-26) & Multiidioma básico; API REST documentada con claves por organización. Los webhooks quedan fuera de la primera entrega. \\
\hline
Trabajo futuro & RF-11 (informe RF-17; revoto, D-03); RF-22 (informe RF-10); RF-23 (informe RF-20); otros elementos aún sin ID local: RF-06, RF-12, RF-13, RF-26 (webhooks) y RF-27 del informe & Facturación, voto preferencial/ponderado, quórum, modo híbrido, revoto, criptografía de extremo a extremo con custodios, webhooks y widget embebible requieren alcance y evaluación adicionales. \\
\hline
\end{longtable}
\normalsize

En esta tabla se usan primero los identificadores locales de la memoria. Los
identificadores entre paréntesis remiten al informe complementario: RF-11
corresponde a RF-17 (revoto), RF-20 a RF-11 (multidioma), RF-21 a RF-26 (API), RF-22 a RF-10 (voto
preferencial) y RF-23 a RF-20 (cifrado en cliente y custodia umbral). Los
requisitos del informe que aún no tienen un requisito local independiente se
identifican expresamente como tales.

La clasificación MoSCoW del informe es una propuesta y no equivale a una
decisión aprobada de implementación. En particular, el voto preferencial y el
cifrado cliente con claves umbral se mantienen como trabajo futuro. En el MVP,
una huella encadenada y un recibo podrán documentar integridad y presencia del
registro, pero no se denominarán verificabilidad criptográfica completa.
También queda pendiente decidir si el canal SMS se incluye o se deja como
adaptador posterior; el correo es el canal mínimo previsto.
